\numoptstrue
\yearpaper{2016}{NEET Phase I \textperiodcentered{} 1 May 2016 \textperiodcentered{} Test Booklet Code A \textperiodcentered{} Physics Q 1--45 \textperiodcentered{} 180 Marks}{}
\begin{multicols}{2}

\QQ{From a disc of radius $R$ and mass $M$, a circular hole of diameter $R$, whose rim passes through the centre is cut. What is the moment of inertia of the remaining part of the disc about a perpendicular axis, passing through the centre?\ntag{X}{}{}{}{}{parallel axis theorem}}
{\opt{$15MR^2/32$}{$13MR^2/32$}{$11MR^2/32$}{$9MR^2/32$}}

\QQ{A square loop ABCD carrying a current $i$, is placed near and coplanar with a long straight conductor XY carrying a current $I$, the net force on the loop will be:
\fig{\begin{tikzpicture}[scale=0.75]
\draw[->] (0,0) -- (0,2.4) node[left,lbl]{Y}; \node[lbl,left] at (0,0) {X}; \draw[->] (-0.2,0.9) -- (-0.2,1.4) node[left,lbl]{$I$};
\draw (0.8,0.3) rectangle (2.4,1.9); \node[lbl,below] at (0.8,0.3) {A}; \node[lbl,above] at (0.8,1.9) {B}; \node[lbl,above] at (2.4,1.9) {C}; \node[lbl,below] at (2.4,0.3) {D};
\draw[->] (1.75,1.3) arc (60:300:0.25); \node[lbl] at (1.75,1.1) {$i$};
\draw[<->] (0,0.0) -- (0.8,0.0) node[midway,below,lbl]{$L/2$}; \draw[<->] (0.8,0.0) -- (2.4,0.0) node[midway,below,lbl]{$L$}; \draw[<->] (2.7,0.3) -- (2.7,1.9) node[midway,right,lbl]{$L$};
\end{tikzpicture}}\ntag{A}{123}{leph104}{17}{ \textperiodcentered{} force between parallel currents}{Class XII Part 1 \textperiodcentered{} Ch 4}}
{\opt{$\dfrac{2\mu_0Ii}{3\pi}$}{$\dfrac{\mu_0Ii}{2\pi}$}{$\dfrac{2\mu_0IiL}{3\pi}$}{$\dfrac{\mu_0IiL}{2\pi}$}}

\QQ{The magnetic susceptibility is negative for:\ntag{D}{147}{leph105}{12}{ \textperiodcentered{} diamagnetic: chi negative}{Class XII Part 1 \textperiodcentered{} Ch 5}}
{\opt{diamagnetic material only}{paramagnetic material only}{ferromagnetic material only}{paramagnetic and ferromagnetic materials}}

\QQ{A siren emitting a sound of frequency 800\,Hz moves away from an observer towards a cliff at a speed of $15\unit{m\,s^{-1}}$. Then, the frequency of sound that the observer hears in the echo reflected from the cliff is: (Take velocity of sound in air $= 330\unit{m\,s^{-1}}$)\ntag{X}{}{}{}{}{Doppler effect}}
{\opt{765\,Hz}{800\,Hz}{838\,Hz}{885\,Hz}}

\QQ{\fig{\begin{circuitikz}[scale=0.62,transform shape,font=\Large]
\draw (0,0) to[battery1,l=$V$] (0,2.4) -- (1.6,2.4) node[circ]{} node[above]{1}; \draw (1.8,2.0) -- (1.6,2.4); \draw (2.6,2.4) node[circ]{} node[above]{2} -- (4,2.4) to[C,l=8\,$\mu$F] (4,0) -- (0,0);
\draw (1.9,1.9) node[right]{S} -- (1.9,1.5) to[C,l_=2\,$\mu$F] (1.9,0);
\end{circuitikz}}
A capacitor of $2\unit{\mu F}$ is charged as shown in the diagram. When the switch S is turned to position 2, the percentage of its stored energy dissipated is:\ntag{A}{76}{leph102}{32}{ \textperiodcentered{} energy shared between capacitors}{Class XII Part 1 \textperiodcentered{} Ch 2}}
{\opt{0\%}{20\%}{75\%}{80\%}}

\QQ{In a diffraction pattern due to a single slit of width `$a$', the first minimum is observed at an angle $30^\circ$ when light of wavelength $5000\unit{\text{\AA}}$ is incident on the slit. The first secondary maximum is observed at an angle of:\ntag{A}{267}{leph202}{13}{ \textperiodcentered{} single slit}{Class XII Part 2 \textperiodcentered{} Ch 10}}
{\opt{$\sin^{-1}\!\left(\dfrac14\right)$}{$\sin^{-1}\!\left(\dfrac23\right)$}{$\sin^{-1}\!\left(\dfrac12\right)$}{$\sin^{-1}\!\left(\dfrac34\right)$}}

\QQ{At what height from the surface of earth the gravitation potential and the value of $g$ are $-5.4\times10^7\unit{J\,kg^{-2}}$ and $6.0\unit{m\,s^{-2}}$ respectively? Take the radius of earth as 6400\,km:\ntag{A}{134}{keph107}{8}{ \textperiodcentered{} V=-GM/r, g=GM/r\textasciicircum{}2}{Class XI Part 1 \textperiodcentered{} Ch 7}}
{\opt{2600\,km}{1600\,km}{1400\,km}{2000\,km}}

\QQ{Out of the following options which one can be used to produce a propagating electromagnetic wave?\ntag{D}{205}{leph108}{5}{ \textperiodcentered{} accelerated charges radiate}{Class XII Part 1 \textperiodcentered{} Ch 8}}
{\opt{A charge moving at constant velocity}{A stationary charge}{A chargeless particle}{An accelerating charge}}

\QQ{Two identical charged spheres suspended from a common point by two massless strings of lengths $l$, are initially at a distance $d$ ($d\ll l$) apart because of their mutual repulsion. The charges begin to leak from both the spheres at a constant rate. As a result, the spheres approach each other with a velocity $v$. Then $v$ varies as a function of the distance $x$ between the spheres, as:\ntag{A}{9}{leph101}{9}{}{Class XII Part 1 \textperiodcentered{} Ch 1}}
{\opt{$v\propto x^{1/2}$}{$v\propto x$}{$v\propto x^{-1/2}$}{$v\propto x^{-1}$}}

\QQ{A uniform rope of length $L$ and mass $m_1$ hangs vertically from a rigid support. A block of mass $m_2$ is attached to the free end of the rope. A transverse pulse of wavelength $\lambda_1$ is produced at the lower end of the rope. The wavelength of the pulse when it reaches the top of the rope is $\lambda_2$. The ratio $\lambda_2/\lambda_1$ is:\ntag{A}{285}{keph207}{8}{ \textperiodcentered{} v=sqrt(T/mu)}{Class XI Part 2 \textperiodcentered{} Ch 14}}
{\opt{$\sqrt{\dfrac{m_1}{m_2}}$}{$\sqrt{\dfrac{m_1+m_2}{m_2}}$}{$\sqrt{\dfrac{m_2}{m_1}}$}{$\sqrt{\dfrac{m_1+m_2}{m_1}}$}}

\QQ{A refrigerator works between $4^\circ$C and $30^\circ$C. It is required to remove 600 calories of heat every second in order to keep the temperature of the refrigerated space constant. The power required is: (Take 1\,cal $= 4.2$\,Joules)\ntag{A}{236}{keph204}{11}{ \textperiodcentered{} refrigerator}{Class XI Part 2 \textperiodcentered{} Ch 11}}
{\opt{2.365\,W}{23.65\,W}{236.5\,W}{2365\,W}}

\QQ{An air column, closed at one end and open at the other, resonates with a tuning fork when the smallest length of the column is $50\unit{cm}$. The next larger length of the column resonating with the same tuning fork is:\ntag{A}{293}{keph207}{16}{ \textperiodcentered{} odd harmonics}{Class XI Part 2 \textperiodcentered{} Ch 14}}
{\opt{66.7\,cm}{100\,cm}{150\,cm}{200\,cm}}

\QQ{Consider the junction diode as ideal. The value of current flowing through AB is:
\fig{\begin{circuitikz}[scale=0.62,transform shape,font=\Large]
\draw (0,0) node[circ]{} node[above]{A} node[below]{$+4$\,V} to[empty diode] (2,0) to[R,l=1\,k$\Omega$] (4.4,0) node[circ]{} node[above]{B} node[below]{$-6$\,V};
\end{circuitikz}}\ntag{A}{334}{leph206}{12}{ \textperiodcentered{} forward bias}{Class XII Part 2 \textperiodcentered{} Ch 14}}
{\opt{0\,A}{$10^{-2}$\,A}{$10^{-1}$\,A}{$10^{-3}$\,A}}

\QQ{The charge flowing through a resistance $R$ varies with time $t$ as $Q = at-bt^2$, where $a$ and $b$ are positive constants. The total heat produced in $R$ is:\ntag{A}{91}{leph103}{11}{ \textperiodcentered{} P=I\textasciicircum{}2R}{Class XII Part 1 \textperiodcentered{} Ch 3}}
{\opt{$\dfrac{a^3R}{6b}$}{$\dfrac{a^3R}{3b}$}{$\dfrac{a^3R}{2b}$}{$\dfrac{a^3R}b$}}

\QQ{A black body is at a temperature of 5760\,K. The energy of radiation emitted by the body at wavelength 250\,nm is $U_1$, at wavelength 500\,nm is $U_2$ and that at 1000\,nm is $U_3$. Wien's constant, $b = 2.88\times10^6\unit{nm\,K}$. Which of the following is correct?\ntag{A}{218}{keph203}{17}{ \textperiodcentered{} Wien's law}{Class XI Part 2 \textperiodcentered{} Ch 10}}
{\opt{$U_1 = 0$}{$U_3 = 0$}{$U_1>U_2$}{$U_2>U_1$}}

\QQ{Coefficient of linear expansion of brass and steel rods are $\alpha_1$ and $\alpha_2$. Lengths of brass and steel rods are $l_1$ and $l_2$ respectively. If $(l_2-l_1)$ is maintained same at all temperatures, which one of the following relations holds good?\ntag{A}{207}{keph203}{6}{ \textperiodcentered{} linear expansion}{Class XI Part 2 \textperiodcentered{} Ch 10}}
{\opt{$\alpha_1l_2 = \alpha_2l_1$}{$\alpha_1l_2^2 = \alpha_2l_1^2$}{$\alpha_1^2l_2 = \alpha_2^2l_1$}{$\alpha_1l_1 = \alpha_2l_2$}}

\QQ{A npn transistor is connected in common emitter configuration in a given amplifier. A load resistance of $800\unit{\Omega}$ is connected in the collector circuit and the voltage drop across it is 0.8\,V. If the current amplification factor is 0.96 and the input resistance of the circuit is $192\unit{\Omega}$, the voltage gain and the power gain of the amplifier will respectively be:\ntag{X}{}{}{}{}{transistor (CE amplifier)}}
{\opt{4, 3.84}{3.69, 3.84}{4, 4}{4, 3.69}}

\QQ{The intensity at the maximum in a Young's double slit experiment is $I_0$. Distance between two slits is $d = 5\lambda$, where $\lambda$ is the wavelength of light used in the experiment. What will be the intensity in front of one of the slits on the screen placed at a distance $D = 10d$?\ntag{A}{266}{leph202}{12}{}{Class XII Part 2 \textperiodcentered{} Ch 10}}
{\opt{$I_0$}{$\dfrac{I_0}4$}{$\dfrac34I_0$}{$\dfrac{I_0}2$}}

\QQ{A uniform circular disc of radius $50\unit{cm}$ at rest is free to turn about an axis which is perpendicular to its plane and passes through its centre. It is subjected to a torque which produces a constant angular acceleration of $2.0\unit{rad\,s^{-2}}$. Its net acceleration in $\mathrm{m\,s^{-2}}$ at the end of 2.0\,s is approximately:\ntag{A}{117}{keph106}{26}{ \textperiodcentered{} rotational kinematics, a\_t and a\_c}{Class XI Part 1 \textperiodcentered{} Ch 6}}
{\opt{8.0}{7.0}{6.0}{3.0}}

\QQ{An electron of mass $m$ and a photon have same energy $E$. The ratio of de-Broglie wavelengths associated with them is: ($c$ being velocity of light)\ntag{A}{285}{leph203}{12}{}{Class XII Part 2 \textperiodcentered{} Ch 11}}
{\opt{$\dfrac1c\left(\dfrac E{2m}\right)^{1/2}$}{$\left(\dfrac E{2m}\right)^{1/2}$}{$c(2mE)^{1/2}$}{$\dfrac1c\left(\dfrac{2m}E\right)^{1/2}$}}

\QQ{A disk and a sphere of same radius but different masses roll off on two inclined planes of the same altitude and length. Which one of the two objects gets to the bottom of the plane first?\ntag{X}{}{}{}{}{kinetic energy of rolling}}
{\opt{Disk}{Sphere}{Both reach at the same time}{Depends on their masses}}

\QQ{The angle of incidence for a ray of light at a refracting surface of a prism is $45^\circ$. The angle of prism is $60^\circ$. If the ray suffers minimum deviation through the prism, the angle of minimum deviation and refractive index of the material of the prism respectively, are:\ntag{A}{240}{leph201}{20}{ \textperiodcentered{} minimum deviation}{Class XII Part 2 \textperiodcentered{} Ch 9}}
{\opt{$45^\circ$; $\dfrac1{\sqrt2}$}{$30^\circ$; $\sqrt2$}{$45^\circ$; $\sqrt2$}{$30^\circ$; $\dfrac1{\sqrt2}$}}

\QQ{When an $\alpha$-particle of mass `$m$' moving with velocity `$v$' bombards on a heavy nucleus of charge `$Ze$', its distance of closest approach from the nucleus depends on $m$ as:\ntag{A}{295}{leph204}{6}{ \textperiodcentered{} distance of closest approach}{Class XII Part 2 \textperiodcentered{} Ch 12}}
{\opt{$\dfrac1m$}{$\dfrac1{\sqrt m}$}{$\dfrac1{m^2}$}{$m$}}

\QQ{A particle of mass 10\,g moves along a circle of radius $6.4\unit{cm}$ with a constant tangential acceleration. What is the magnitude of this acceleration if the kinetic energy of the particle becomes equal to $8\times10^{-4}\unit{J}$ by the end of the second revolution after the beginning of the motion?\ntag{A}{42}{keph103}{16}{}{Class XI Part 1 \textperiodcentered{} Ch 3}}
{\opt{$0.1\unit{m/s^2}$}{$0.15\unit{m/s^2}$}{$0.18\unit{m/s^2}$}{$0.2\unit{m/s^2}$}}

\QQ{The molecules of a given mass of a gas have r.m.s.\ velocity of $200\unit{m\,s^{-1}}$ at $27^\circ$C and $1.0\times10^5\unit{N\,m^{-2}}$ pressure. When the temperature and pressure of the gas are respectively, $127^\circ$C and $0.05\times10^5\unit{N\,m^{-2}}$, the r.m.s.\ velocity of its molecules in $\mathrm{m\,s^{-1}}$ is:\ntag{A}{252}{keph205}{9}{ \textperiodcentered{} v\_rms proportional to sqrt(T)}{Class XI Part 2 \textperiodcentered{} Ch 12}}
{\opt{$100\sqrt2$}{$\dfrac{400}{\sqrt3}$}{$\dfrac{100\sqrt2}3$}{$\dfrac{100}3$}}

\QQ{A long straight wire of radius $a$ carries a steady current $I$. The current is uniformly distributed over its cross-section. The ratio of the magnetic fields $B$ and $B'$, at radial distances $\dfrac a2$ and $2a$ respectively, from the axis of the wire is:\ntag{D}{119}{leph104}{13}{ \textperiodcentered{} Example 4.7}{Class XII Part 1 \textperiodcentered{} Ch 4}}
{\opt{$\dfrac14$}{$\dfrac12$}{1}{4}}

\QQ{A particle moves so that its position vector is given by $\vec r = \cos\omega t\,\hat x+\sin\omega t\,\hat y$, where $\omega$ is a constant. Which of the following is true?\ntag{A}{42}{keph103}{16}{ \textperiodcentered{} uniform circular motion}{Class XI Part 1 \textperiodcentered{} Ch 3}}
{\opt{Velocity and acceleration both are perpendicular to $\vec r$.}{Velocity and acceleration both are parallel to $\vec r$.}{Velocity is perpendicular to $\vec r$ and acceleration is directed towards the origin.}{Velocity is perpendicular to $\vec r$ and acceleration is directed away from the origin.}}

\QQ{What is the minimum velocity with which a body of mass $m$ must enter a vertical loop of radius $R$ so that it can complete the loop?\ntag{D}{79}{keph105}{9}{ \textperiodcentered{} Example 5.7}{Class XI Part 1 \textperiodcentered{} Ch 5}}
{\opt{$\sqrt{gR}$}{$\sqrt{2gR}$}{$\sqrt{3gR}$}{$\sqrt{5gR}$}}

\QQ{When a metallic surface is illuminated with radiation of wavelength $\lambda$, the stopping potential is $V$. If the same surface is illuminated with radiation of wavelength $2\lambda$, the stopping potential is $\dfrac V4$. The threshold wavelength for the metallic surface is:\ntag{A}{278}{leph203}{5}{}{Class XII Part 2 \textperiodcentered{} Ch 11}}
{\opt{$4\lambda$}{$5\lambda$}{$\dfrac52\lambda$}{$3\lambda$}}

\QQ{A gas is compressed isothermally to half its initial volume. The same gas is compressed separately through an adiabatic process until its volume is again reduced to half. Then:\ntag{A}{233}{keph204}{8}{ \textperiodcentered{} adiabatic curve is steeper}{Class XI Part 2 \textperiodcentered{} Ch 11}}
{\opt{Compressing the gas isothermally will require more work to be done.}{Compressing the gas through adiabatic process will require more work to be done.}{Compressing the gas isothermally or adiabatically will require the same amount of work.}{Which of the case (whether compression through isothermal or through adiabatic process) requires more work will depend upon the atomicity of the gas.}}

\QQ{A potentiometer wire is $100\unit{cm}$ long and a constant potential difference is maintained across it. Two cells are connected in series first to support one another and then in opposite direction. The balance points are obtained at $50\unit{cm}$ and $10\unit{cm}$ from the positive end of the wire in the two cases. The ratio of emf's is:\ntag{X}{}{}{}{}{potentiometer}}
{\opt{5:1}{5:4}{3:4}{3:2}}

\QQ{A astronomical telescope has objective and eyepiece of focal lengths $40\unit{cm}$ and $4\unit{cm}$ respectively. To view an object $200\unit{cm}$ away from the objective, the lenses must be separated by a distance:\ntag{A}{243}{leph201}{23}{}{Class XII Part 2 \textperiodcentered{} Ch 9}}
{\opt{37.3\,cm}{46.0\,cm}{50.0\,cm}{54.0\,cm}}

\QQ{Two non-mixing liquids of densities $\rho$ and $n\rho$ ($n>1$) are put in a container. The height of each liquid is $h$. A solid cylinder of length $L$ and density $d$ is put in this container. The cylinder floats with its axis vertical and length $pL$ ($p<1$) in the denser liquid. The density $d$ is equal to:\ntag{A}{184}{keph202}{5}{ \textperiodcentered{} upthrust from both liquids}{Class XI Part 2 \textperiodcentered{} Ch 9}}
{\opt{$\{1+(n+1)p\}\rho$}{$\{2+(n+1)p\}\rho$}{$\{2+(n-1)p\}\rho$}{$\{1+(n-1)p\}\rho$}}

\QQ{To get output 1 for the following circuit, the correct choice for the input is:
\fig{\begin{circuitikz}[scale=0.7,transform shape]
\draw (0,0.6) node[or port](G1){}; \draw (2.4,0.1) node[and port](G2){};
\draw (G1.in 1) -- ++(-0.5,0) node[ocirc]{} node[left]{A}; \draw (G1.in 2) -- ++(-0.5,0) node[ocirc]{} node[left]{B};
\draw (G1.out) -| (G2.in 1); \draw (G2.in 2) -- ++(-3.2,0) node[ocirc]{} node[left]{C}; \draw (G2.out) -- ++(0.5,0) node[ocirc]{} node[below]{Y};
\end{circuitikz}}\ntag{X}{}{}{}{}{digital electronics (logic gates)}}
{\opt{A = 0, B = 1, C = 0}{A = 1, B = 0, C = 0}{A = 1, B = 1, C = 0}{A = 1, B = 0, C = 1}}

\QQ{A piece of ice falls from a height $h$ so that it melts completely. Only one-quarter of the heat produced is absorbed by the ice and all energy of ice gets converted into heat during its fall. The value of $h$ is: [Latent heat of ice is $3.4\times10^5\unit{J/kg}$ and $g = 10\unit{N/kg}$]\ntag{A}{212}{keph203}{11}{ \textperiodcentered{} latent heat}{Class XI Part 2 \textperiodcentered{} Ch 10}}
{\opt{34\,km}{544\,km}{136\,km}{68\,km}}

\QQ{The ratio of escape velocity at earth ($v_e$) to the escape velocity at a planet ($v_p$) whose radius and mean density are twice as that of earth is:\ntag{A}{136}{keph107}{10}{ \textperiodcentered{} v\_e=sqrt(2GM/R)}{Class XI Part 1 \textperiodcentered{} Ch 7}}
{\opt{1:2}{$1:2\sqrt2$}{1:4}{$1:\sqrt2$}}

\QQ{If the magnitude of sum of two vectors is equal to the magnitude of difference of the two vectors, the angle between these vectors is:\ntag{A}{33}{keph103}{7}{}{Class XI Part 1 \textperiodcentered{} Ch 3}}
{\opt{$0^\circ$}{$90^\circ$}{$45^\circ$}{$180^\circ$}}

\QQ{Given the value of Rydberg constant is $10^7\unit{m^{-1}}$, the wave number of the last line of the Balmer series in hydrogen spectrum will be:\ntag{A}{291}{leph204}{2}{ \textperiodcentered{} Balmer formula}{Class XII Part 2 \textperiodcentered{} Ch 12}}
{\opt{$0.025\times10^4\unit{m^{-1}}$}{$0.5\times10^7\unit{m^{-1}}$}{$0.25\times10^7\unit{m^{-1}}$}{$2.5\times10^7\unit{m^{-1}}$}}

\QQ{A body of mass 1\,kg begins to move under the action of a time dependent force $\vec F = (2t\hat i+3t^2\hat j)\unit{N}$, where $\hat i$ and $\hat j$ are unit vectors along $x$ and $y$ axis. What power will be developed by the force at the time $t$?\ntag{A}{83}{keph105}{13}{ \textperiodcentered{} P=F.v}{Class XI Part 1 \textperiodcentered{} Ch 5}}
{\opt{$(2t^2+3t^3)$\,W}{$(2t^2+4t^4)$\,W}{$(2t^3+3t^4)$\,W}{$(2t^3+3t^5)$\,W}}

\QQ{An inductor $20\unit{mH}$, a capacitor $50\unit{\mu F}$ and a resistor $40\unit{\Omega}$ are connected in series across a source of emf $V = 10\sin340t$. The power loss in A.C.\ circuit is:\ntag{A}{186}{leph107}{10}{}{Class XII Part 1 \textperiodcentered{} Ch 7}}
{\opt{0.51\,W}{0.67\,W}{0.76\,W}{0.89\,W}}

\QQ{If the velocity of a particle is $v = At+Bt^2$, where $A$ and $B$ are constants, then the distance travelled by it between 1\,s and 2\,s is:\ntag{A}{18}{keph102}{6}{ \textperiodcentered{} calculus method}{Class XI Part 1 \textperiodcentered{} Ch 2}}
{\opt{$\dfrac32A+4B$}{$3A+7B$}{$\dfrac32A+\dfrac73B$}{$\dfrac A2+\dfrac B3$}}

\QQ{A long solenoid has 1000 turns. When a current of 4\,A flows through it, the magnetic flux linked with each turn of the solenoid is $4\times10^{-3}\unit{Wb}$. The self-inductance of the solenoid is:\ntag{D}{167}{leph106}{14}{ \textperiodcentered{} N Phi = L I}{Class XII Part 1 \textperiodcentered{} Ch 6}}
{\opt{4\,H}{3\,H}{2\,H}{1\,H}}

\QQ{A small signal voltage $V(t) = V_0\sin\omega t$ is applied across an ideal capacitor C:\ntag{D}{183}{leph107}{7}{ \textperiodcentered{} average power in a capacitor is zero}{Class XII Part 1 \textperiodcentered{} Ch 7}}
{\opt{Current $I(t)$, lags voltage $V(t)$ by $90^\circ$.}{Over a full cycle the capacitor C does not consume any energy from the voltage source.}{Current $I(t)$ is in phase with voltage $V(t)$.}{Current $I(t)$ leads voltage $V(t)$ by $180^\circ$.}}

\QQ{Match the corresponding entries of column 1 with column 2. [Where $m$ is the magnification produced by the mirror]\par\smallskip
\begin{tabular}{@{}ll@{\qquad}ll@{}}
\multicolumn{2}{@{}l}{\textit{Column 1}} & \multicolumn{2}{l}{\textit{Column 2}}\\
(A) & $m = -2$ & (a) & Convex mirror\\
(B) & $m = -\frac12$ & (b) & Concave mirror\\
(C) & $m = +2$ & (c) & Real image\\
(D) & $m = +\frac12$ & (d) & Virtual image
\end{tabular}\ntag{A}{226}{leph201}{6}{ \textperiodcentered{} sign of magnification}{Class XII Part 2 \textperiodcentered{} Ch 9}}
{\opt{A$\to$b and c; B$\to$b and c; C$\to$b and d; D$\to$a and d}{A$\to$a and c; B$\to$a and d; C$\to$a and b; D$\to$c and d}{A$\to$a and d; B$\to$b and c; C$\to$b and d; D$\to$b and c}{A$\to$c and d; B$\to$b and d; C$\to$b and c; D$\to$a and d}}

\QQ{A car is negotiating a curved road of radius $R$. The road is banked at an angle $\theta$. The coefficient of friction between the tyres of the car and the road is $\mu_s$. The maximum safe velocity on this road is:\ntag{D}{64}{keph104}{16}{ \textperiodcentered{} banked road}{Class XI Part 1 \textperiodcentered{} Ch 4}}
{\opt{$\sqrt{gR^2\dfrac{\mu_s+\tan\theta}{1-\mu_s\tan\theta}}$}{$\sqrt{gR\dfrac{\mu_s+\tan\theta}{1-\mu_s\tan\theta}}$}{$\sqrt{\dfrac gR\dfrac{\mu_s+\tan\theta}{1-\mu_s\tan\theta}}$}{$\sqrt{\dfrac g{R^2}\dfrac{\mu_s+\tan\theta}{1-\mu_s\tan\theta}}$}}

%%BEGIN-ANALYSIS
\end{multicols}
\par\needspace{16\baselineskip}\vspace{14pt}
\begin{center}{\pgfornament[width=3.2cm,color=gilt]{87}}\par
{\large\color{navy}\scshape NEET 2016 \textperiodcentered{} Answer Key}\end{center}
{\small\setlength{\tabcolsep}{1pt}\renewcommand{\arraystretch}{1.5}
\begin{tabularx}{\linewidth}{@{}*{10}{>{\centering\arraybackslash}X}@{}}\toprule
\textcolor{gray}{\textbf{1}\,(2)$^{\dagger}$} & \textbf{2}\,(1) & \textbf{3}\,(1) & \textcolor{gray}{\textbf{4}\,(3)$^{\dagger}$} & \textbf{5}\,(4) & \textbf{6}\,(4) & \textbf{7}\,(1) & \textbf{8}\,(4) & \textbf{9}\,(3) & \textbf{10}\,(2)\\
\textbf{11}\,(3) & \textbf{12}\,(3) & \textbf{13}\,(2) & \textbf{14}\,(1) & \textbf{15}\,(4) & \textbf{16}\,(4) & \textcolor{gray}{\textbf{17}\,(1)$^{\dagger}$} & \textbf{18}\,(4) & \textbf{19}\,(1) & \textbf{20}\,(1)\\
\textcolor{gray}{\textbf{21}\,(2)$^{\dagger}$} & \textbf{22}\,(2) & \textbf{23}\,(1) & \textbf{24}\,(1) & \textbf{25}\,(2) & \textbf{26}\,(3) & \textbf{27}\,(3) & \textbf{28}\,(4) & \textbf{29}\,(4) & \textbf{30}\,(2)\\
\textcolor{gray}{\textbf{31}\,(4)$^{\dagger}$} & \textbf{32}\,(4) & \textbf{33}\,(4) & \textcolor{gray}{\textbf{34}\,(4)$^{\dagger}$} & \textbf{35}\,(3) & \textbf{36}\,(2) & \textbf{37}\,(2) & \textbf{38}\,(3) & \textbf{39}\,(4) & \textbf{40}\,(1)$^{*}$\\
\textbf{41}\,(3) & \textbf{42}\,(4) & \textbf{43}\,(2) & \textbf{44}\,(1) & \textbf{45}\,(2) &  &  &  &  & \\\bottomrule\end{tabularx}}\par\smallskip
{\footnotesize Answers solved by the editor (the official final key of this year could not be obtained from the NTA/CBSE websites). (1)--(4) are the printed options. {\color{gray}Grey$^{\dagger}$}: deleted portion (answer given for completeness). $^{*}$: see note.\par
\par\smallskip\textbf{Q40.} The exact value is about 0.46\,W; (1) 0.51\,W is the nearest option and is the intended answer.
}\par
\clearpage
\begin{center}{\pgfornament[width=5cm,color=gilt]{87}}\par
{\fontsize{20}{24}\selectfont\color{navy}\scshape NEET 2016 \textperiodcentered{} Paper Analysis}\par
{\color{gilt}\pgfornament[width=5cm,color=gilt,symmetry=h]{87}}\end{center}
\vspace{4pt}
\begin{multicols}{2}
\textbf{\color{crimson}1. Direct vs Applied}\par\smallskip
{\footnotesize Checked question by question against the four 2026-27 NCERT books (XI Part 1 \& 2, XII Part 1 \& 2). Out of 45, 39 questions are in the present syllabus; percentages are of these.}\par\smallskip
\begin{tabular}{@{}lrr@{}}\toprule
Direct from NCERT & \textbf{7} & 18\%\\
\quad text, summary, table or figure & 5 &\\
\quad NCERT solved Example & 2 &\\
\quad NCERT Exercise & 0 &\\
Applied (concept in NCERT) & \textbf{32} & 82\%\\\midrule
Analysed & 39 &\\
Deleted portion (not analysed) & 6 &\\\midrule
\textbf{Total} & \textbf{45} &\\\bottomrule\end{tabular}\par\smallskip
{\footnotesize Direct: Q 3, 8, 26, 28, 42, 43, 45.\\ From NCERT Examples: Q 26 (Example 4.7), 28 (Example 5.7).}\par\medskip
\textbf{\color{crimson}2. Difficulty level}\par\smallskip
\begin{tabular}{@{}lr@{}}\toprule
Easy & 19\\ Moderate & 17\\ Hard & 3\\\bottomrule\end{tabular}\par\smallskip
{\footnotesize Hard: Q 9, 14, 18. The 2016 paper (Phase I) was moderate to hard: many questions combine two ideas, and six were on topics now deleted. Class XI: 19 questions, Class XII: 20 (analysed questions only). (Difficulty is the editor's assessment.)}\par\medskip
\columnbreak
\textbf{\color{crimson}3. Chapter-wise distribution}\par\smallskip
{\footnotesize\begin{tabular}{@{}p{4.7cm}c@{\,}r@{\,}r@{\,}r@{}}\toprule
Chapter & Cl. & Qs & Dir & App\\\midrule
Motion in a Straight Line & XI & 1 & 0 & 1\\
Motion in a Plane & XI & 3 & 0 & 3\\
Laws of Motion & XI & 1 & 1 & 0\\
Work, Energy \& Power & XI & 2 & 1 & 1\\
Rotational Motion & XI & 1 & 0 & 1\\
Gravitation & XI & 2 & 0 & 2\\
Mech.\ Prop.\ of Fluids & XI & 1 & 0 & 1\\
Thermal Prop.\ of Matter & XI & 3 & 0 & 3\\
Thermodynamics & XI & 2 & 0 & 2\\
Kinetic Theory & XI & 1 & 0 & 1\\
Waves & XI & 2 & 0 & 2\\
Electric Charges \& Fields & XII & 1 & 0 & 1\\
Electrostatic Potential \& Capacitance & XII & 1 & 0 & 1\\
Current Electricity & XII & 1 & 0 & 1\\
Moving Charges \& Magnetism & XII & 2 & 1 & 1\\
Magnetism \& Matter & XII & 1 & 1 & 0\\
Electromagnetic Induction & XII & 1 & 1 & 0\\
Alternating Current & XII & 2 & 1 & 1\\
Electromagnetic Waves & XII & 1 & 1 & 0\\
Ray Optics & XII & 3 & 0 & 3\\
Wave Optics & XII & 2 & 0 & 2\\
Dual Nature & XII & 2 & 0 & 2\\
Atoms & XII & 2 & 0 & 2\\
Semiconductor Electronics & XII & 1 & 0 & 1\\
\bottomrule\end{tabular}}\par\medskip
\textbf{\color{crimson}4. Teaching takeaways}\par\smallskip
{\footnotesize\begin{itemize}[leftmargin=1em,nosep]
\item Examples 4.7 (field inside and outside a thick wire) and 5.7 (vertical loop, $\sqrt{5gR}$) gave direct marks.
\item Six questions were on deleted portions (parallel axis theorem, Doppler effect, transistor amplifier, rolling, potentiometer, logic gates).
\item 32 of the 39 in-syllabus questions are applied; the hard ones (leaking charges, $Q=at-bt^2$, YDSE in front of a slit) need a short derivation.
\item Circular motion was asked four times (tangential acceleration, $\vec r=\cos\omega t\,\hat x+\sin\omega t\,\hat y$, vertical loop, banked road).
\end{itemize}}\par\medskip
\textbf{\color{crimson}5. Deleted portions}\par\smallskip
\textbf{6} of the 45 questions are on topics removed from the 2026-27 NCERT (rationalised syllabus). They are marked \textsc{Deleted Portion} and are not analysed: parallel axis theorem (Q 1); Doppler effect (Q 4); transistor (CE amplifier) (Q 17); kinetic energy of rolling (Q 21); potentiometer (Q 31); digital electronics (logic gates) (Q 34).\par
\end{multicols}
%%END-ANALYSIS


